%% kapa: exam=1 round=27 year=2016 period=1 form=A subject=economics subject_ko=경제학원론 seq=14 number=54
%% source: Q-Net 감정평가사 기출문제 (cst00309) artlSeq=5169402 file=extracted/1008/5169402/제27회 감정평가사 1차 문제/제27회 감정평가사 1교시문제(A형).pdf
%% conversion: kordoc 4.24.51 -> tools/gen1.py (auto); review=proofread-sonnet; flags=none
\begin{question}{54}
이윤극대화를 추구하는 독점기업과 완전경쟁기업의 차이점에 관한 설명으로 옳지 않은 것은?
\begin{choices}
  \item 독점기업의 한계수입은 가격보다 낮은 반면, 완전경쟁기업의 한계수입은 시장가격과 같다.
  \item 독점기업의 한계수입곡선은 우상향하는 반면, 완전경쟁기업의 한계수입곡선은 우하향한다.
  \item 독점기업이 직면하는 수요곡선은 우하향하는 반면, 완전경쟁기업이 직면하는 수요곡선은 수평이다.
  \item 단기균형에서 독점기업은 가격이 한계비용보다 높은 점에서 생산하는 반면, 완전경쟁기업은 시장가격과 한계비용이 같은 점에서 생산한다.
  \item 장기균형에서 독점기업은 경제적 이윤을 얻을 수 있는 반면, 완전경쟁기업은 경제적 이윤을 얻을 수 없다.
\end{choices}
\end{question}
\answer{2}
